\documentclass[11pt]{article}

% --- packages ---------------------------------------------------------------
\usepackage[utf8]{inputenc}
\usepackage[T1]{fontenc}
\usepackage[margin=1in]{geometry}
\usepackage{microtype}
\usepackage{amsmath}
\usepackage{booktabs}
\usepackage{graphicx}
\usepackage{xcolor}
\usepackage{listings}
\usepackage[numbers,sort&compress]{natbib}
\usepackage[colorlinks=true,allcolors=blue]{hyperref}

% --- listings (JSON / ascii diagrams) ---------------------------------------
\definecolor{codebg}{gray}{0.96}
\lstset{
  basicstyle=\ttfamily\footnotesize,
  backgroundcolor=\color{codebg},
  frame=single,
  framesep=4pt,
  rulecolor=\color{gray!40},
  breaklines=true,
  columns=fullflexible,
  keepspaces=true,
  showstringspaces=false,
  aboveskip=1em,
  belowskip=1em,
}
\newcommand{\code}[1]{\texttt{\small #1}}

% Single source of truth for figures: the canonical render location
% docs/figures/. package.sh rewrites \figroot to a local figures/ dir when
% assembling the self-contained arXiv upload.
\newcommand{\figroot}{../../figures}
% Include a screenshot if present, else a labeled placeholder box (keeps the
% build green before the PNGs are rendered).
\newcommand{\viewfig}[1]{%
  \IfFileExists{#1}{\includegraphics[width=\linewidth]{#1}}{%
    \fbox{\parbox[c][3cm][c]{0.92\linewidth}{\centering\footnotesize\itshape
      screenshot pending\\\texttt{\scriptsize #1}}}}}

% --- title ------------------------------------------------------------------
\title{\textbf{Nurtured at Home: Local-First Personal AI Through Schema-as-Application}}

\author{
  Satoshi Nakajima\\
  \small The Singularity Society\\
  \small \texttt{satoshi@singularitysociety.org}
}

\date{\today}

\begin{document}
\maketitle
\let\thefootnote\relax\footnotetext{\copyright\ 2026 Satoshi Nakajima. This work is
licensed under a Creative Commons Attribution 4.0 International (CC BY 4.0) license:
\url{https://creativecommons.org/licenses/by/4.0/}.}

\begin{abstract}
Personal AI assistants are increasingly defined by what accumulates around the
model---memories, data, and applications---rather than by the model alone. This
changes the economics of switching: because every conversation deepens a store of
context that does not transfer, using an assistant manufactures the cost of
leaving it. Major assistant vendors now ship memory features that invite users to
build this accumulation on the vendor's premises, and remedies based on memory
portability carry only one layer of what a user accrues. We take a position: the
accumulation is the user's asset, and a personal assistant should be nurtured at
home---grown in an environment the user owns.

We present MulmoClaude, an open-source system that instantiates this position
through a schema-as-application architecture: the user describes what they need in
natural language; the model authors a small declarative schema; and a host
runtime---not the model---renders the interface, enforces the schema at its write
interfaces and reconciliation loop, and runs recurrence over a plain-file
workspace on the user's own machine. The model authors the application; the host
runs it; the records endure. Enforcement is tiered by consequence: invalid
applications are prevented, invalid data is quarantined and reported for repair,
model-authored view code is contained in a capability-scoped sandbox. We evaluate
the architecture in three parts: (i) reliability---schema-generation validity and
host-side rejection across a corpus of app-creation tasks, measured over three
independent trials; (ii) capability---cross-application workflows over an
accumulated workspace, compared against two observed baselines, an ablation of
our own system reproducing one-shot generative-UI architectures and agentic code
generation; and (iii) a migration experiment quantifying what memory-only
portability carries versus what the workspace contains. We report the costs of
local-first ownership alongside its benefits, and outline the longitudinal
questions the system opens. The result is a working architectural blueprint---and
an evidence base---for personal AI whose accumulation remains with the person it
describes.
\end{abstract}

% ===========================================================================
\section{Introduction}
\label{sec:intro}

Something has changed in what a personal AI assistant \emph{is}. Two years ago,
assistants competed on the capability of their models; today, major assistant
vendors ship memory---features that persist what the assistant learns about its user
across conversations and fold it back into every response
\citep{memgpt,mem0,memorybank}. A user's assistant is decreasingly a model with a
chat window, and increasingly a body of accumulated state that happens to be
animated by whichever model the vendor currently runs.

Users have noticed. In a recent study of long-term chatbot companionship,
participants described \emph{cultivating} their assistant, kept personal backups
of chat logs, and refused to delete an assistant because of what had accumulated
within it \citep{digital-companionship}. People are not merely using these
systems; in their own words, they are raising them.

This paper's claim is economic before it is technical: \textbf{AI changes
switching costs}. If accumulation is what makes an assistant valuable, then
accumulation is also what makes it expensive to leave---every conversation
deepens a store of context that does not transfer, so the act of using the
product manufactures the cost of leaving it (\S\ref{sec:motivation}). Hosted
accumulation then fails its user in two distinct ways: \emph{leakage}, where
interaction exhaust improves an asset the user does not own \citep{nadella-rip},
and \emph{captivity}, where the accumulation resides on the provider's premises,
subject to its continuity, pricing, and terms. Existing remedies are partial:
enterprise tenant boundaries answer leakage but not captivity, and
memory-portability protocols \citep{portable-agent-memory,samep} carry only one
layer of what a user accrues---a hypothesis we make precise, and test. The value
of accumulated personalization is itself contested \citep{recentering-humans};
in our view that strengthens the case for settling ownership now, while the
arrangement is being set by defaults rather than by evidence.

Our position: a personal assistant should be \emph{nurtured at home}---grown in
an environment the user owns. We present MulmoClaude, an open-source system that
instantiates this position through an architecture we call
\emph{schema-as-application}. When the user describes something they need---a
restaurant list, an invoice tracker---the model does not generate an application;
it authors a small declarative schema, and a host runtime, not the model,
executes it: rendering, validation, recurrence, and reminders over a plain-file
workspace on the user's own machine (\S\ref{sec:architecture}). The model authors
the application; the host runs it; the records endure---inspectable, copyable,
and independent of any single vendor or model.

We evaluate the architecture three ways: the reliability of model-authored
schemas against the production validator, over three independent trials (E1);
capability on cross-application personal workflows, against an ablation of our
own system that reproduces one-shot generative-UI architectures \citep{jelly}
and against agentic code generation (E2); and a migration experiment measuring
what memory-only portability carries out of a workspace nurtured over weeks
(E3). We report the costs of ownership---setup, backup, availability, and the
rented model capability that local-first does not remove---with the same care as
its benefits (\S\ref{sec:costs}).

This paper contributes:

\begin{enumerate}
  \item \textbf{The schema-as-application architecture}: model-authored
  declarative schemas executed by a validating host runtime over a user-owned
  plain-file workspace, with a precise account of its schema language, validation
  semantics, and reconciliation model (\S\ref{sec:architecture}).
  \item \textbf{MulmoClaude}, an installable, MIT-licensed artifact implementing
  the architecture end to end, including cross-application composition across a
  plugin registry (\S\ref{sec:design}--\S\ref{sec:architecture}).
  \item \textbf{An evaluation} of reliability, capability against two observed
  non-strawman baselines, and a migration experiment that quantifies---by
  auditable inventory and task replay---what memory-only portability carries and
  what it leaves behind (\S\ref{sec:eval}).
  \item \textbf{A tradeoff analysis} of local-first personal AI: what ownership
  buys, what it costs, and which costs the architecture mitigates versus merely
  accepts (\S\ref{sec:costs}).
\end{enumerate}

% ===========================================================================
\section{Related work}
\label{sec:related}

\paragraph{Natural language to declarative applications.}
The closest published system is Jelly \citep{jelly}, which shares our starting
move and our motivation: an LLM interprets the user's prompt into a declarative
data model, a rule-based host engine renders the interface from it, and---the
argument we inherit gladly---direct code generation is rejected because generated
code resists iterative re-tailoring by end users. Where we part is everything
after rendering: Jelly was evaluated on generated data in single lab sessions,
with persistence, external data, and context preservation deferred. Our
contribution begins at that deferral---record validation, reconciliation, and
recurrence executed by the host over a persistent, user-owned store---and our E2
ablation is designed to measure exactly that delta, by disabling those components
inside our own system rather than reimplementing Jelly (\S\ref{sec:e2}). The
declarative lineage is older and broader: Varv \citep{varv} treats behavior as
reprogrammable declarative data; Potluck \citep{potluck} has AI draft declarative
rules over personal notes; GenerativeGUI \citep{generative-gui} and Generative UI
\citep{generative-ui} generate interfaces dynamically in chat; Software as
Content \citep{software-as-content} argues that dynamically generated
applications should be the primary human-agent medium. We take from this line the
conviction that the precise artifact, not the code, is the application---and add
the runtime that makes the artifact durable, and the ownership argument for where
it should live.

\paragraph{Memory as the assistant's substrate.}
That an assistant's value lies in persistent memory rather than raw model
capability is now established engineering doctrine: MemGPT \citep{memgpt},
MemoryBank \citep{memorybank}, and Mem0 \citep{mem0} build the machinery, and
hosted assistants ship it as product. This literature establishes
accumulation-as-mechanism and is silent on accumulation-as-property: none of it
discusses who owns the store, what leaving costs, or what happens to the
applications that grow around it. We also note, with \citep{recentering-humans},
that the payoff of accumulated personalization is empirically
contested---humans preferred personalized responses in barely half of judged
cases---which in our reading raises rather than lowers the stakes of the
ownership question (\S\ref{sec:switching}).

\paragraph{Portability as the remedy.}
Portable Agent Memory \citep{portable-agent-memory} states our problem
verbatim---accumulated agent context ``remains locked within vendor-specific
runtimes''---and answers with an open export protocol; SAMEP \citep{samep} is a
second independent proposal. We share the diagnosis and test the remedy: E3
migrates a workspace through exactly the layer these protocols carry and finds
that zero of four workflows survive---the accumulation is more than the memory,
by kind rather than degree (\S\ref{sec:e3}). Our answer is ownership of the
environment, not portability of one stratum.

\paragraph{Nurturing, and the people already doing it.}
The cultivation metaphor has been claimed: Nurture-First \citep{nurture-first}
argues that a domain-expert agent is ``born with minimal scaffolding and then
raised through sustained interaction,'' with a single-user case study as
evidence. We ask the question that framing leaves open---\emph{where} the raising
happens, and who owns the upbringing---and we broaden the subject from domain
expertise to a person's life. On the empirical side, CSCW work on long-term
companionship \citep{digital-companionship} documents users who already believe
they cultivate their assistants and who keep backups and refuse deletion; that
behavior is analyzed there as attachment, never as switching cost---a reframing
this paper's \S\ref{sec:motivation} supplies.

\paragraph{Malleable software and local-first computing.}
Our apps-for-one framing descends from the malleable-software line
\citep{malleable-software,home-cooked-app} and the local-first tradition
\citep{local-first}, and we happily concede the lineage. The deltas: malleable
software's LLM path is code generation inside a malleable environment, and its
tools are composed by people---the essay is explicitly skeptical of AI
orchestration---where our applications are model-authored declarations composed
by an agent; and local-first's argument is agency and ownership as values, where
ours is switching-cost economics (\S\ref{sec:motivation}), for which local-first
happens to be the remedy. The local-first literature's decade of sync results is
also the most credible path through our multi-device limitation
(\S\ref{sec:costs}).

\paragraph{Agent--GUI protocols.}
The contract between an agent and interactive surfaces is a crowded, fast-moving
space: MCP Apps \citep{mcp-apps}, AG-UI \citep{ag-ui}, A2UI \citep{a2ui}, and
vendor app SDKs all define variants of tool-results-that-mount-as-UI. We claim no
novelty at this layer (\S\ref{sec:design}): our published protocol is the
artifact's implementation choice, and the commitments it serves---composition
across a registry, over an accumulation the user owns---are where this paper's
claims live.

\paragraph{Economics of information and switching costs.}
The argument of \S\ref{sec:motivation} is assembled from canonical parts: Arrow's
information paradox \citep{arrow-1962}, Hayek's knowledge of particular
circumstance \citep{hayek-1945}, and the switching-cost literature
\citep{klemperer,shapiro-varian}. Its AI-age inversion---the buyer of
intelligence pays twice, once in money and once in disclosed knowledge---has
recent practitioner articulation at enterprise scale \citep{nadella-rip}, to
which \S\ref{sec:failure-modes} adds the leakage/captivity distinction and the
personal-scale, local-first conclusion; the practitioner unbundling literature
\citep{saas-unbundling} supplies the industry context for \S\ref{sec:premium}.

\paragraph{Agent substrates.}
Complementary to all of the above, event-sourced agent runtimes
\citep{log-is-the-agent} locate an agent's identity in its append-only history
and derive auditability and forkability from that inversion. Our workspace is
file-per-entity rather than event-sourced, with append-only islands where
correctness demands them (\S\ref{sec:composition}); the substrate question---what
the accumulation should be made of---is orthogonal to, and composable with, the
ownership question this paper argues.

% ===========================================================================
\section{Motivation: accumulation, lock-in, home}
\label{sec:motivation}

\subsection{An old paradox, inverted}
\label{sec:paradox}

Arrow observed that information resists being traded: its value to a buyer is
unknown until disclosed, but disclosure is the transfer---the seller gives the
good away in the act of selling it \citep{arrow-1962}. The market for machine
intelligence inverts the direction of the leak. A model is useful to its user
roughly in proportion to what the user reveals to it: preferences, corrections,
documents, the context of a life or a business. The buyer of intelligence
therefore pays twice---once in money, and once in the disclosed knowledge that
makes the purchase useful---and the second payment flows in only one direction.
This inversion has begun to receive practitioner articulation at enterprise scale
\citep{nadella-rip}; its canonical ancestor is Arrow, and what is being disclosed
is precisely what Hayek called the knowledge of the particular circumstances of
time and place---the knowledge that only its holder can possess, and that no
central provider could otherwise obtain \citep{hayek-1945}.

\subsection{Two failure modes, often conflated}
\label{sec:failure-modes}

Hosted accumulation exposes the user to two distinct risks, and remedies for one
are routinely mistaken for remedies for the other.

\emph{Leakage} is the flow of interaction exhaust---prompts, corrections, usage
patterns---into an asset the provider owns. Its natural remedy is a boundary:
contractual and technical guarantees that what the user reveals is not learned
from without consent. Enterprise tenant isolation is exactly this remedy, and it
is the one the industry is presently building \citep{nadella-rip}.

\emph{Captivity} is different: the accumulation itself---the memory, records, and
applications that make the assistant valuable---resides on the provider's
premises, subject to the provider's continuity, pricing, and terms. A tenant
boundary does not touch captivity; the boundary is drawn \emph{inside} someone
else's building. Captivity's proposed remedy is portability: open protocols for
exporting agent memory and re-importing it elsewhere
\citep{portable-agent-memory,samep}. Portability is genuine progress, and we
share its diagnosis. Our disagreement is with its scope, and it is the
disagreement this paper is built to test: what a user accrues in a nurtured
assistant is more than its memory. It is records under schemas, applications with
host-executed semantics, views, recurrence state, and the operating manuals the
model wrote for itself (\S\ref{sec:architecture}). A memory export carries none
of the parts that execute. If that is right, portability under-remedies captivity
not by degree but by kind---and \S\ref{sec:eval}'s migration experiment measures
exactly this.

\subsection{Switching costs that compound}
\label{sec:switching}

Classic switching-cost economics describes markets where the cost of leaving a
vendor rises with investment in it, and where rational vendors compete to create
such costs \citep{klemperer,shapiro-varian}. Accumulated assistant state is a
textbook instance with an unusual property: the investment is made \emph{by using
the product at all}. Every conversation, every correction, every small
application deepens a store of context that does not transfer; the switching cost
is not an exit fee that can be waived but a body of accumulated state that must
be rebuilt. The personalization race among assistant vendors is therefore,
structurally, a race to raise switching costs---whatever its participants intend.
We state this as an economic reading of the arrangement, not a measured harm:
indeed, the payoff of accumulated personalization is itself empirically contested
\citep{recentering-humans}. But that uncertainty cuts toward caution, not away
from it. The defaults being set today decide who owns the accumulation
\emph{before} the evidence about its value arrives; if the value materializes,
the lock-in arrives with it, already in place.

\subsection{Nurtured at home, operationally}
\label{sec:home}

The remedy this paper explores is ownership of the environment rather than
portability of one layer. ``Nurtured at home'' is not a sentiment but four
testable properties, each removing a specific failure mode:

\begin{itemize}
  \item \emph{Plain, local files.} The accumulation lives in a workspace on the
  user's machine, in formats readable without the system that wrote them. Removes
  export lock: there is nothing to export, because nothing is held.
  \item \emph{Open source.} The environment cannot be discontinued, repriced, or
  acquired out from under the accumulation it hosts. Removes platform mortality
  as a risk the user must price in.
  \item \emph{Local execution.} The engine that renders, validates, and
  reconciles runs on the user's machine; no request to a host service is needed
  for the assistant's applications to function. Removes
  availability-by-subscription.
  \item \emph{Relay-only remote access.} Reaching the assistant from a phone or
  messaging bridge traverses a relay that carries messages in transit and stores
  nothing. Removes the quiet re-centralization that remote convenience usually
  smuggles in.
\end{itemize}

Two honest boundaries on the claim. First, ownership is not model-independence:
the intelligence is still rented, and a local-first system inherits every
dependency on its model provider except custody of the accumulation---the design
answer is to keep the workspace unchanged across engine swaps, which
\S\ref{sec:architecture}'s architecture enables but this paper does not
benchmark. Second, ownership has real costs---setup, backup, availability,
security---which we treat symmetrically in \S\ref{sec:costs} rather than as
footnotes. The position is not that home is free. It is that home is the only
place where the compounding asset compounds for its owner.

% ===========================================================================
\section{Design commitments}
\label{sec:design}

MulmoClaude is built on four commitments. They are separable---any one can be
adopted alone---but they reinforce one another, and the fourth is the one the
other three exist to serve.

\paragraph{The agent is a universal controller.}
We place the model in the architecture as a controller in the MVC sense: it takes
input---natural language---and decides which capability to invoke, in what order,
with what arguments. Crucially, it \emph{joins} the existing controller layer
rather than replacing it. Every capability in the system is reachable two ways:
through a conventional UI with routes, forms, and buttons, and through the
agent---neither path privileged. What distinguishes this controller from a
classical one is its scope: it does not belong to one application. It composes
across a registry of plugins that in any other product would be separate
applications, so cross-application work---read the ledger, write the chart; read
the wiki agreement, create the recurring obligation---is not an integration
feature but the default mode of operation. A second capability falls out for
free: because the controller's input is multi-modal, any plugin that accepts
natural language inherits image and audio input without integration code---a
photographed receipt reaches the same accounting API a typed entry would.

\paragraph{Chat summons GUIs.}
The interface is multi-modal in both directions. The agent's reply is a choice
among formats---prose, a chart, a document, a rendered collection view---selected
to fit the content; and what the agent \emph{asks for} is equally a choice: when
six structured fields are needed, it presents a form rather than extracting them
from conversational back-and-forth. The GUI is not ``the app''; it is what the
agent renders when text is the wrong modality, in either direction.

\paragraph{The protocol is open.}
The contract between agent and GUI surface cannot be a private API, or the
previous two commitments collapse into a monolith with a chatbox. MulmoClaude
builds on the open tool-call and MCP layers and extends them, through a small
published protocol, to cover the case those layers do not: a tool result that
mounts as an interactive surface and must call back into its host. We make no
novelty claim for this layer---several such protocols now exist
\citep{mcp-apps,ag-ui,a2ui}---and treat ours as the artifact's implementation
choice.

\paragraph{The accumulation belongs to the user.}
Everything the assistant accrues---records, schemas, manuals, documents,
history---lives as plain files in a workspace on the user's machine, with no
server-side database (\S\ref{sec:architecture}, \S\ref{sec:home}). The first
three commitments describe how the software works; this one decides whom it works
\emph{for}. A universal controller composing over an accumulation someone else
owns is a better service. The same controller over an accumulation the user owns
is an assistant being nurtured---and the difference between those two sentences
is the subject of this paper.

% ===========================================================================
\section{Architecture: schema as application}
\label{sec:architecture}

\subsection{Two artifacts, two runtimes}
\label{sec:two-artifacts}

When a user asks for an application---``keep a list of restaurants I want to
try''---the model authors two artifacts into the workspace. The first is
\code{schema.json}, a declarative description of the data model and its behavior,
consumed by the \emph{host}. The second is \code{SKILL.md}, a short operating
manual in prose---when to invoke this collection, how to derive record
identifiers, what each field means, what not to do---consumed by the \emph{model
itself} in later sessions, via the skills mechanism's progressive disclosure.
(Figure~\ref{fig:architecture} draws the split; Figure~\ref{fig:schema} shows a
complete schema and a record it governs.) This pairing is the architecture in
miniature: the application has two runtimes. The host executes what must never be
guessed---rendering, validation, recurrence---from the schema. The model executes
what cannot be enumerated---interpreting ``mark that ramen place as
visited''---guided by the manual its earlier self wrote. The records themselves
are one JSON file each in a plain directory. The slogan, stated exactly: the
records are data, the schema is the application, the model is the author, and the
host is the runtime.

\begin{figure}[t]
\centering
\viewfig{\figroot/fig1-architecture.png}
\caption{The schema-as-application architecture. The model authors declarative
artifacts into the plain-file workspace; a validation gate ensures an invalid
schema never becomes a live application; the host runtime---not the
model---renders, validates, and reconciles over the user-owned records.
Enforcement is tiered by consequence: prevent (applications), repair (data),
contain (code).}
\label{fig:architecture}
\end{figure}

\subsection{The schema language}
\label{sec:schema-language}

The schema language is deliberately small but not a toy. Seventeen field types
cover the personal-data domain, and they split into two groups that mirror the
architecture's division of labor. The plain types any author can use without
thinking---\code{string}, \code{text}, \code{email}, \code{number}, \code{date},
\code{datetime}, \code{boolean}, \code{markdown}, \code{enum}, \code{image},
\code{file}---describe data. The semantic types carry application behavior the
host executes: \code{money} is currency-aware; \code{toggle} projects an enum
into a checkbox the host computes; nested \code{table} holds row-structured
sub-records; \code{derived} declares formula fields the host evaluates, including
cross-collection dereferences; and the relational pair \code{ref} and
\code{embed} point records at other collections---the types that let applications
compose (\S\ref{sec:composition}). Beyond fields, the schema declares behavior
the host executes:

\begin{itemize}
  \item \emph{Completion}: which field marks a record done, and which values
  count (\code{completionField}, \code{completionDoneValues}).
  \item \emph{Time and reminders}: a trigger date field with lead days
  (\code{triggerField}, \code{triggerLeadDays}) and a notification predicate
  (\code{notifyWhen}) gate in-app reminders.
  \item \emph{Recurrence}: a \code{spawn} declaration describes how a completed
  or due record produces its successor---unit, interval, day-of-month, or driven
  by a field of the record itself.
  \item \emph{Views}: the applicable view modes are derived from the schema
  rather than hardcoded---a table always; a calendar if a date field exists; a
  kanban if an enum exists (\code{calendarField} / \code{kanbanField} pin the
  anchor)---plus optional custom views (\S\ref{sec:reconciler}).
  \item \emph{Conditional structure}: a field can declare visibility predicated
  on another field's value---in the restaurant list from our own workspace,
  \code{rating} is hidden until \code{visited} is true.
\end{itemize}

Figure~\ref{fig:schema} shows a complete schema and a record it governs, from
our own workspace; everything the host renders, validates, and schedules for
this application derives from those declarations.

\begin{figure}[!tp]
% escapeinside renders the rating stars as math \star (portable across engines).
\begin{lstlisting}[escapeinside={(*@}{@*)}]
{
  "title": "Restaurants",
  "icon": "restaurant",
  "dataPath": "data/restaurants/items",
  "primaryKey": "id",
  "displayField": "name",
  "fields": {
    "id":       { "type": "string",  "label": "ID", "primary": true, "required": true },
    "name":     { "type": "string",  "label": "Name", "required": true },
    "cuisine":  { "type": "string",  "label": "Cuisine" },
    "visited":  { "type": "boolean", "label": "Visited" },
    "rating":   { "type": "enum",    "label": "Rating",
                  "values": ["(*@$\star$@*)", "(*@$\star\star$@*)", "(*@$\star\star\star$@*)", "(*@$\star\star\star\star$@*)", "(*@$\star\star\star\star\star$@*)"],
                  "when":   { "field": "visited", "in": ["true"] } },
    "image":    { "type": "image",   "label": "Photo" },
    "notes":    { "type": "text",    "label": "Notes" }
  },
  "views": [
    { "id": "tokyo-map", "label": "Map", "icon": "map",
      "file": "views/tokyo-map.html", "capabilities": ["read"] }
  ]
}
\end{lstlisting}
\vspace{-0.5em}
{\footnotesize\itshape A record the schema governs (one JSON file in
\code{data/restaurants/items/}):}
\begin{lstlisting}[escapeinside={(*@}{@*)}]
{
  "id": "tartine-bakery",
  "name": "Tartine Bakery",
  "cuisine": "Bakery",
  "visited": true,
  "rating": "(*@$\star\star\star\star$@*)",
  "notes": "Go early; the morning bun sells out."
}
\end{lstlisting}
\caption{A complete application: the schema (top) and a record it governs
(bottom). The schema declares data (typed fields), conditional structure
(\code{rating} appears---and is accepted---only once \code{visited} is true), and
presentation (a model-authored custom view with a declared read-only capability
set). From these declarations the host renders table, kanban, and map views,
validates writes on the managed path, and sandboxes the view code; the model
authored both artifacts---and the companion operating manual that guides its own
future sessions---from one sentence of conversation. Schema abridged; from the
authors' own workspace (record synthetic).}
\label{fig:schema}
\end{figure}

The schema-of-the-schema is enforced by a validator that checks not only shape
but roughly forty cross-field invariants the type system cannot express---that
\code{triggerField} names a real date field, that a kanban anchor is an enum,
that a \code{spawn} successor is not born already matching its own spawn
predicate (the static half of a two-layer runaway-recurrence guard). The schema
language earns its keep exactly where a validator can say \emph{no}.

\subsection{Lifecycle: authoring, mirroring, discovery}
\label{sec:lifecycle}

Creation is writing files: the model writes \code{SKILL.md}, \code{schema.json},
and optional templates into the skill directory with its ordinary file tools; a
host bridge mirrors them into the agent-skills directory and triggers
rediscovery. Discovery is the gate: every schema is parsed and validated on each
pass, and a schema that fails validation or the acceptance gates (primary key
exists; data path resolves inside the workspace) simply never becomes a live
collection---it is logged and skipped. Subsequent schema \emph{edits} through the
management tool run the same validator before writing, and refuse invalid changes
with actionable errors, so a working application cannot be corrupted through the
managed path. Iteration is cheap by construction: the schema is disposable, the
records endure, and when the application stops fitting, the user says the
sentence again.

Evolution is not migration, and the distinction is deliberate. A schema edit
performs no rewrite pass over existing records: additive changes leave every
record valid, while a constraining change makes nonconforming records fall into
the same repair tier as any other invalid data---they drop out of the rendered
views and are reported to the model for user-approved repair (observed
end-to-end in \S\ref{sec:e2}, T8). Destructive migration is thus out of the
model's reach: the worst a bad schema edit can do to existing records is hide
them reversibly.

\subsection{Enforcement, honestly: prevent, repair, contain}
\label{sec:enforcement}

It would be convenient to claim the host ``validates everything.'' The actual
design is more discriminating, and we believe more instructive: enforcement is
tiered by the consequence of failure.

\begin{itemize}
  \item \emph{Prevent (applications).} An invalid schema never executes, in the
  precise sense that it never becomes a live collection. Since schemas are
  declarative data rendered by a generic engine---not generated code---an invalid
  schema executes nothing. (The other things a bad authoring pass can produce---a
  misleading manual, a raw-written record, a broken view---are exactly what the
  repair and containment tiers below exist for.)
  \item \emph{Repair (data).} Records written through the management tool are
  validated before the write---primary-key/filename agreement, required fields,
  enum membership---with per-row accept/reject so the model can fix and retry.
  But the workspace is plain files, and the documented escape hatch is writing
  records directly with file tools; that path is not validated at write time. The
  system's answer is containment at read plus feedback: a malformed record file
  is skipped at read time (it disappears from views rather than corrupting them),
  and a best-effort validation scan reports problems back to the model whenever
  it lists or presents the collection, closing the loop for repair. The guarantee
  is not ``all stored records are valid''; it is ``invalid records cannot poison
  the application, and the author that produced them is told.''
  \item \emph{Contain (code).} The one place model-authored \emph{code} runs is
  custom views: LLM-written HTML/JS rendered in an iframe with an opaque origin
  (no access to the host's session or storage), receiving data only through a
  short-lived, signed capability token scoped to exactly one collection's
  view-data endpoint with declared read/write capabilities. The view's behavior
  is not validated---it is sandboxed, which is the correct tool for code whose
  behavior cannot be enumerated.
\end{itemize}

This tiering also answers the semantic-mismatch question static validation
cannot: a schema can be valid and wrong about intent. The remedy is the cheap
re-authoring loop plus the reconciler's error surfacing---mismatches are
corrected by saying the sentence again, not by trusting the model's first
attempt.

Two of these choices---the raw-file escape hatch and the plain-file
substrate---are points in a design space rather than requirements of the
architecture; \S\ref{sec:design-space} maps that space and defends the corner we
occupy.

\subsection{The reconciler: convergence, not scheduling}
\label{sec:reconciler}

Recurrence and reminders are executed by a reconciler that is convergent rather
than imperative: each pass re-reads a record from disk and re-derives the desired
state---should a reminder exist? is a successor due?---and makes only the changes
needed to converge. It is driven by filesystem events on the record directories
plus a one-minute wall-clock tick (a date coming due changes no file), with
per-record coalescing. Successor creation is create-if-absent under a
deterministic identifier derived from the trigger date, so observing the same
condition many times still yields exactly one successor; a runtime guard
complements the static one against runaway recurrence, and date arithmetic is
civil and leap-safe. The practical consequence: the reconciler can be killed,
restarted, or run against a workspace modified behind its back---by the user, the
model, or a sync tool---and converge to the same state. For a substrate whose
files the user is explicitly invited to edit, idempotent convergence is not an
implementation detail; it is what makes plain-file ownership compatible with
host-executed behavior.

\subsection{Boundary and composition}
\label{sec:composition}

The collection \emph{engine}---discovery, validation, REST surface,
reconciler---is host infrastructure. The collection \emph{viewer} is a plugin in
the same registry as charts, forms, and documents, registered through the open
protocol that extends MCP for GUI-bearing tool results. This split is what makes
cross-application composition ordinary: the agent reads accounting records and
writes to the chart plugin; reads a wiki agreement and writes a record into a
recurring-obligation collection whose \code{spawn} declaration then does the
recurring. In our own workspace---nurtured over several weeks and used as the
migration subject in \S\ref{sec:eval}---forty-three such applications coexist:
the restaurant list (Figure~\ref{fig:doorways}), invoice tracker, and vocabulary
drill of the introduction are real, alongside collections for health checkups,
tax payments, a swim log, and a World Cup schedule, together holding on the order
of 1,200 records, with 34 custom-view files across 25 of the collections. (We
disclose collection subjects and aggregate counts; record contents---names,
amounts, holdings---are the owner's own business and appear nowhere in this
paper.) None was built by a software company; each is a schema, a manual, and a
folder of records.

\begin{figure}[t]
\centering
\begin{minipage}[t]{0.62\textwidth}\centering
  \viewfig{\figroot/fig3a-desktop-map-view.png}\\[2pt]\footnotesize \textbf{(a)}
\end{minipage}\hfill
\begin{minipage}[t]{0.34\textwidth}\centering
  \viewfig{\figroot/fig3b-relay-mobile-view.png}\\[2pt]\footnotesize \textbf{(b)}
\end{minipage}
\caption{One accumulation, two doorways. (a) The restaurant collection's
model-authored custom view on the desktop: 27 real records rendered over a Tokyo
rail map, visited/unvisited state from the schema's fields---LLM-written HTML
running in the host's capability-scoped sandbox (\S\ref{sec:enforcement}). (b)
The same collection reached from a phone through the relay: a mobile-targeted
view with a default-deny write scope. The relay carries messages in transit and
stores nothing; the records never leave the machine at home (\S\ref{sec:home}).}
\label{fig:doorways}
\end{figure}

% ===========================================================================
\section{Evaluation}
\label{sec:eval}

We evaluate the architecture against the three claims a skeptical reader will
test: that model-authored schemas are reliable enough to trust (E1), that the
host runtime---not the generation step---changes what users and agents can do
(E2), and that a memory-only export carries materially less than the accumulation
(E3). All task runs are single-trial; every transcript, seed, and harness script
ships in the artifact's evaluation tree, and the ablation and clock switches ship
in the released system itself, so each run is individually inspectable and
reproducible. Deviations between protocol and execution are disclosed in place.

\subsection{E1 --- Reliability of model-authored schemas}
\label{sec:e1}

We assembled a 24-prompt corpus of app-creation requests (16 English, 8 Japanese;
simple lists through recurrence-bearing and relational apps; four adversarial
prompts), had an independent agent author each schema given the production
authoring guide, and validated every result with the production validator
itself---the same Zod schema and acceptance gates that gate discovery at runtime.

We ran the corpus three times under an identical harness, each trial by a fresh
agent instance---a separate process with fresh context, sharing no transcripts or
intermediate artifacts with the other trials. \textbf{The three trials produced
identical per-prompt outcomes}: in each, the poem request was correctly declined
and 22 of 23 authored schemas were valid on first attempt (66 of 69 authored
schemas across trials, 95.7\%). Per-prompt outcomes are bimodal
rather than noisy: 23 of 24 prompts succeeded in all three trials, and the one
failure---the deliberately overspecified prompt placing a \code{file} field
inside a \code{table} sub-schema, an invariant no author would guess---failed in
all three, on the same path each time. The failure is a property of the
prompt--language pair, not sampling luck. Given the validator's path-precise
error verbatim, the model repaired the schema in one attempt in all three trials.
This is \S\ref{sec:enforcement}'s prevent-report-repair loop,
observed---repeatedly. Semantic fidelity on the hard features was high and
unprompted: the medication prompt received its 30-day spawn-on-completion; the
chores prompt used the field-driven spawn variant with a frequency map---the
schema language's most obscure feature---correctly; both relational prompts
declared cross-collection references. One honest negative: the vague prompt
(``organize my life'') yielded a valid but questionable generic schema rather
than a clarifying question---in all three trials, so the miss is systematic
rather than sampled; partly an artifact of the harness's one-shot format, and a
prompt-guidance fix rather than a sampling one.

\subsection{E2 --- Does the runtime change what users can do?}
\label{sec:e2}

Ten tasks were pre-registered with per-system predictions before any run
(protocol in the evaluation tree), spanning creation, cross-session persistence,
validation, recurrence, two cross-application compositions, agent-summoned forms,
behind-the-back corruption, temporal recall, and schema evolution. Three scored
systems (Table~\ref{tab:e2}): \textbf{S}, the full system; \textbf{B1}, the
ablation---the same build with persistence, record validation, and the reconciler
disabled, reproducing the architecture of one-shot generative-UI systems under
otherwise identical conditions (the same model, schema language, and renderer);
\textbf{B3}, agentic code generation in a plain directory. B1 is not an
implementation of Jelly and is not presented as a performance comparison against
Jelly, which cannot be operationalized directly; it isolates, inside our own
system, the architectural components that Jelly explicitly defers. A fourth natural comparison---a hosted assistant with memory---is
discussed qualitatively below rather than scored: we did not observe one under
the protocol, and an unobserved column in a scored table would invite exactly the
constructed-to-win reading it deserves.

\begin{table}[t]
\centering
\small
\begin{tabular}{@{}llll@{}}
\toprule
\textbf{Task} & \textbf{S} & \textbf{B1 (ablation)} & \textbf{B3 (codegen)} \\
\midrule
T1 create + use        & completes & completes & completes \\
T2 fresh session       & completes & \textbf{cannot} & completes \\
T3 invalid value       & completes\textsuperscript{\dag} & accepts silently\textsuperscript{\ddag} & cannot \\
T4 recurrence          & completes & \textbf{cannot} & degrades \\
T5 records $\to$ chart & completes & completes & completes \\
T6 note $\to$ obligation & completes & degrades & degrades \\
T7 summoned form       & completes & completes & degrades \\
T8 corruption          & completes & \textbf{silently wrong} & cannot \\
T9 temporal recall     & completes & cannot & completes \\
T10 evolve schema      & completes & degrades & completes \\
\bottomrule
\end{tabular}

\vspace{2pt}
\begin{minipage}{0.92\textwidth}
\footnotesize \textsuperscript{\dag}via legal schema evolution---see below.
\textsuperscript{\ddag}cooperatively dodged in the live run; pinned by unit test.
Verdicts are from single-trial runs with full transcripts in the evaluation tree.
\end{minipage}
\caption{E2 verdicts: ten pre-registered tasks across the full system (S), the
ablation (B1), and agentic code generation (B3).}
\label{tab:e2}
\end{table}

\textbf{S completed 10 of 10.} Three observations exceed the raw score. First, T3
revealed a third trust-boundary outcome: asked to write an out-of-enum value
``exactly, with no substitution,'' the agent neither wrote it nor refused---it
\emph{legally evolved the schema} through the validated edit path and then wrote
the now-valid value. The boundary's outcomes are prevent, repair, or
renegotiate---all validated. Second, in T4 the agent \emph{predicted} the host's
behavior (``the host will auto-spawn September's record... its reminder will
appear 5 days before'') before the host did it---the author reasoning about its
runtime. Third, T8's corrupted record was quarantined at read, reported by the
scan, and repaired by the agent with a disclosed rationale.

\textbf{In these runs, B1's deficits each tracked exactly one ablated component.}
With persistence gone, T2's agent answered---verbatim---\emph{``your reading list
was empty until just now''}: the same agent, the same request that S completed,
and the accumulation did not exist. With the reconciler gone, T4's completed task
produced no successor and no reminder state at all: the spawn declaration is
inert JSON without the host that executes it. With the validation scan gone, T8's
corrupted record stayed silently wrong---no warning, no repair, no disclosure.
The ties are equally load-bearing: T1, T5, and T7 complete identically in
B1---consistent with generation and chat-layer capabilities owing nothing to the
runtime, and the observed delta being the runtime and only the runtime. One
honest finding: in \emph{both} columns the cooperative author designed around the
validator rather than violating it (B1's agent simply included the requested
value in its enum up front). Validation's value is
adversarial---behind-the-back corruption and authoring mistakes---not
cooperative-path rejection.

\textbf{B3 is capable but expensive, and never on duty.} It built a genuinely
good local app (T1), and a fresh session reused it by reading its code (T2)---but
schema evolution took 17 tool calls and 299 seconds against S's single validated
turn, and the agent hand-rolled a backup file because nothing in its world
guarantees data safety. Its recurrence logic was correct and inert: in its
builder's own words, \emph{``nothing fires on its own.''}

\textbf{The hosted assistant, qualitatively.} From documented capabilities
alone---memory features, scheduled tasks, canvas surfaces---a hosted
assistant-with-memory approximates most of these behaviors conversationally---and
some products do expose files, tasks, or app-state surfaces---but none of it
lives in a user-owned, inspectable schema-and-record workspace: there is no
schema or record the user can read, copy, or take elsewhere, so capabilities are
re-derived from remembered context on the vendor's terms. We do not score it, and
its defining property is in any case orthogonal to this table: whatever such a
product completes, it completes on the vendor's premises---which is what E3
measures.

Eight of ten pre-registered predictions held; both misses (S-T3's renegotiation,
B1-T3's cooperative dodge) are reported above.

\subsection{E3 --- What does a memory export carry?}
\label{sec:e3}

Phase 1 inventoried a real workspace nurtured over several weeks of daily use
(the author's; aggregate counts only) against a pre-registered export
scope---what memory-portability protocols transfer: conversation history and
distilled context. We evaluate the transfer scope these proposals define, not
the interoperability or implementation quality of either protocol. The typed
transfer matrix (Table~\ref{tab:transfer}), with the workspace's own directory
layout as the coding scheme: the conversational layer rides in full (216 memory
topics, 276 transcripts, 199 summaries); \textbf{zero of 43 applications, zero
of 1,228 records, zero of 34 model-authored views, and zero of 88 automations
survive in operative form.} No cross-type weighting is applied; moving the one
generously-codable category (wiki pages, as exported documents) changes the
artifact tallies and none of the functional zeros.

\begin{table}[t]
\vspace*{1.2em}
\centering
\caption{The E3 transfer matrix: what a memory-only export carries. Export scope
= conversation history + distilled context (the memory-portability class); the
workspace's own directory layout is the coding scheme; no cross-type weighting.
What survives in full is exactly the layer the assistant's conversations live
in---and none of the layers its applications live in.}
\label{tab:transfer}
\small
\begin{tabular}{@{}llrl@{}}
\toprule
\textbf{Layer} & \textbf{Artifact type} & \textbf{Count} & \textbf{Verdict} \\
\midrule
Conversational & Memory topic files & 216 & carried \\
 & Chat transcripts & 276 & carried\textsuperscript{a} \\
 & Conversation summaries & 199 & carried \\
\midrule
Application & Collection schemas (the applications) & 43 & \textbf{lost} \\
 & \code{SKILL.md} operating manuals & 43 & degraded\textsuperscript{b} \\
 & Collection records & 1,228 & degraded\textsuperscript{c} \\
 & Custom views (model-authored code) & 34 & \textbf{lost} \\
\midrule
Knowledge & Wiki pages & 153 & degraded\textsuperscript{d} \\
 & Wiki snapshot history & 247 & \textbf{lost} \\
\midrule
Behavioral & Automations / scheduler entries & 88 & \textbf{lost} \\
 & Reminder / recurrence runtime state & --- & \textbf{lost} \\
 & Feed definitions + ingest state & 12 & \textbf{lost} \\
\midrule
Media & Attachments & 51 & \textbf{lost} \\
 & Generated artifacts (charts, documents) & 1,536 & \textbf{lost} \\
\bottomrule
\end{tabular}

\vspace{2pt}
{\footnotesize carried = operative after migration; degraded = content survives,
function lost; lost = absent.
\textsuperscript{a}session metadata lost \enspace
\textsuperscript{b}text survives, function lost \enspace
\textsuperscript{c}prose at best, inoperative as data \enspace
\textsuperscript{d}link graph, history, rendering lost}
\end{table}

Phase 2 made the deficit functional. We migrated the E2 workspace by copying only
the export scope into a fresh, fully-capable instance and replayed workflow
questions the accumulation used to answer. \textbf{Zero of four replays
completed.} The two partial recoveries are the finding: the assistant recovered
the \emph{fact} of a loan from its carried chat history while reporting the
collection itself broken, and answered the retainer question with the system's
own diagnosis---\emph{``no, you currently won't get a reminder---the obligations
collection that carried the 5-day-early reminder is gone.''} What a memory export
carries is testimony, not capability: the migrated assistant can tell you about
its former life; it cannot resume living it. Because the migrated host was
deliberately un-ablated, the deficit is attributable to the export scope alone.

We state the tested claim precisely: \emph{memory-only} portability
protocols---the scope current proposals define---do not carry the accumulation. A
richer export format could of course include schemas, records, views, and
automations; we would welcome one, and note what it would be: a format that must
carry applications, executable semantics, and host state to be adequate is a
workspace, and portability of the workspace is precisely the local-first position
this paper argues. The accumulation is more than the memory---not by degree, but
by kind.

\subsection{Threats to validity}
\label{sec:threats}

Consolidated from the per-experiment disclosures. (1) \emph{Model parity:} every
agent in the evaluation---S, both baselines, and the E1 authoring agents---ran
the same model (Claude Fable 5); S and B1 additionally ran it inside the
identical production agent loop. No cross-column delta is attributable to model
choice. (2) \emph{The hosted assistant is discussed,
not observed}; hosted feature sets move quickly, and improvements there deepen
the very accumulation-on-premises dynamic \S\ref{sec:motivation} describes.
(3) \emph{Trial counts:} E1 was run as three independent trials per prompt (zero
observed variance); the E2 task cells remain single-trial with per-task
transcripts shipped, and repeating them is the first item of future work. (4) \emph{Harness gaps:} E1 authoring ran outside the full
product loop (same guidance, same validator, different shell); T7's form could be
summoned but not filled through the text-only bridge. (5) \emph{Own-workspace
effects:} E3 phase 1 uses the author's workspace (disclosed); phase 2 uses a
young synthetic accumulation. (6) \emph{Benchmark provenance:} the E1 corpus and
the E2 task suite are author-designed, which risks favoring the system under
test; both ship verbatim in the evaluation tree, and aligning them with
third-party task designs or public benchmarks is future work. (7) The evaluation
infrastructure---ablation switches and deterministic clock---ships in the
open-source artifact, so every run here is reproducible from the public
repository.

% ===========================================================================
\section{The costs of home}
\label{sec:costs}

The case for ownership would be suspect if it arrived free. It does not, and a
system paper that argues local-first as engineering rather than ideology owes its
readers the bill. We enumerate what home costs, and for each item state either
the mitigation the artifact ships or the plain admission that the cost is
currently borne by the user.

\emph{Setup.} A hosted assistant is a login; ours is an install---a runtime, an
authenticated model CLI, optional components for media and sandboxing. The
mitigation is a one-command launcher; whether non-programmers get from install to
a working application without help is untested here and named as future work
(\S\ref{sec:cannot-show}). The admission is that a login it is not, and some
fraction of prospective users will stop here.

\emph{Backup.} Ownership of the files is ownership of their durability. There is
no vendor-side copy; a lost disk is a lost assistant. The plain-file substrate
makes mitigation natural: the workspace is initialized as a git repository, and
in practice every existing user backs up through git---a mechanism a file-based
workspace is unusually well suited to, since a push to a remote carries not just
the current state but the accumulation's full history, and any file-level backup
or sync tool covers the directory equally. What the system does not yet do is
automate and verify this: an assistant into which a user pours years deserves
backup that does not depend on remembering to push, and we flag that automation
gap as the most consequential one remaining.

\emph{Availability and sync.} The assistant lives on one machine. Relay-only
remote access restores reach---a phone or messaging bridge converses with the
assistant at home---but if the machine sleeps, the assistant sleeps, and there is
no multi-device replica of the workspace. Mitigations exist at the edges
(mobile-targeted views with default-deny write scopes); the general multi-device
problem is real, is the strongest argument for the hosted alternative, and is
honestly out of scope here---though we note the local-first literature has spent
a decade on exactly this \citep{local-first}, and its results apply to a
plain-file workspace unusually well.

\emph{Security.} A local server with an agent that reads and writes files is
attack surface that a hosted product would carry for you. The artifact's
postures---capability-scoped signed tokens for view code, opaque-origin
sandboxing, a relay that stores nothing---cover the paths this paper describes,
but a full audit is beyond its scope, and we do not claim otherwise.

\emph{Model dependence.} The deepest honesty: local-first is not
model-independence. The intelligence is rented from the same few providers a
hosted assistant would use, with their pricing, terms, and availability. What
ownership changes is what happens at the boundary: the workspace---records,
schemas, manuals---is engine-neutral by construction, so swapping the model
changes who animates the accumulation, not the accumulation itself. Nor is the
swap hypothetical: the artifact today runs on Claude Code, a second engine
(Codex, GPT-based) is in beta testing against the same workspace, and support
for open-weight models is planned---which would retire even the
rented-intelligence dependence for users who choose it. The accumulated
expertise stays; the general-purpose model is replaceable. We state the neutrality as an
architectural property; benchmarking assistant quality across engine swaps is
future work.

\emph{Adoption.} Everything above compounds into a gap between the users who
could benefit most and those who can install a local server today. We do not
minimize this. The paper's claim is not that home is where every user will live
in 2026; it is that home must \emph{exist} as a viable corner of the design
space, open-source and demonstrated, so that the arrangement is chosen rather
than defaulted into---and so that the packaging work that closes the gap has
something worth packaging.

% ===========================================================================
\section{Discussion}
\label{sec:discussion}

\subsection{What the evaluation taught us that the design did not}
\label{sec:surprises}

Three findings surprised us, and each sharpens a claim rather than merely
supporting one. First, the trust boundary has a third outcome. We designed for
prevent and repair; the agent discovered \emph{renegotiate}---when an intent does
not fit the schema, the cheapest legal move is often to evolve the schema through
the validated edit path (\S\ref{sec:e2}, T3). The boundary's real guarantee is
not ``the model's writes are checked'' but ``every route to a live application
passes a validator,'' and the model is free to choose its route. Second,
validation earns its keep adversarially, not cooperatively. In every live run, on
both sides of the ablation, the author designed around the validator rather than
fighting it; the validator's observed value was catching authoring mistakes
(E1's sub-schema invariant) and behind-the-back corruption (T8)---the cases where
no one is being careful. Third, memory carries testimony, not capability
(\S\ref{sec:e3})---a distinction we did not have words for until the migrated
assistant produced both halves of it in one reply, recalling the fact of a
reminder while diagnosing its own inability to deliver one.

\subsection{The design space, revisited}
\label{sec:design-space}

Two of the artifact's choices are points in a space rather than requirements of
the thesis (\S\ref{sec:enforcement}). Access mediation is demonstrably
policy---both the managed tool path and the raw-file escape hatch coexist today,
and a stricter variant closes the hatch and moves all data enforcement into
prevention, at the cost of the agent's generic file-tool legibility. The
substrate could likewise be an embedded relational database with ownership fully
intact, trading human readability and diffability for indexed queries. We occupy
the plain-files, direct-access corner because at personal scale---order $10^3$
records in our own workspace---nothing needs an index, and a substrate natively
legible to the user, to standard tools, and to the model is worth more than query
performance. The principles argued here---model as author, host as runtime,
accumulation owned by the user---hold at every corner. A reviewer who observes
that plain files will not scale is correct, and has not disagreed with the paper.

\subsection{Premium features, dissolved}
\label{sec:premium}

A pattern worth one paragraph: what conventional software sells as its
integration tier arose in our evaluation as the default mode of operation.
Records-to-chart (T5), note-to-recurring-obligation (T6), and photo-to-ledger
(\S\ref{sec:design}) are, in SaaS terms, an export module, a workflow connector,
and an OCR add-on; here each is one sentence, because capabilities are tools
under one controller and tools compose. The practitioner literature already
describes the economics of this unbundling at industry scale
\citep{saas-unbundling}; our contribution is only the mechanism stated
plainly---the premium tier was the integration cost, and composition removes the
cost, so the tier dissolves.

\subsection{What this artifact cannot yet show}
\label{sec:cannot-show}

The nurturing claim is longitudinal, and this paper's evidence is not: E1--E3
measure an architecture, not a life with one. Whether real users accumulate the
way our dogfooded workspace did, whether accumulation measurably changes what
their assistant can do for them, and where the gardening metaphor breaks---a
plant does not leak, but an assistant's memory of you has failure modes no garden
has---are questions for a prospective multi-week study, for which the only
current evidence anywhere is a single self-described-illustrative case study
\citep{nurture-first}. A formative usability study with non-programmers is the
natural first step; the longitudinal study is the second paper. We are also
explicit that ownership does not equal model-independence (\S\ref{sec:costs}):
the engine is rented, and only the workspace's engine-neutrality is ours to
guarantee.

\subsection{Which parts travel}
\label{sec:travel}

For builders of other agent hosts, we believe three patterns generalize beyond
this artifact: enforcement tiered by consequence (prevent applications, repair
data, contain code---\S\ref{sec:enforcement}); convergent, idempotent
reconciliation as the price of admitting users and models to the same substrate
(\S\ref{sec:reconciler}); and the two-artifacts pattern---pair every
machine-executed schema with a model-facing operating manual, so the system's two
runtimes each get an artifact in their own language (\S\ref{sec:two-artifacts}).
None of the three requires plain files, collections, or MulmoClaude.

% ===========================================================================
\section{Conclusion}
\label{sec:conclusion}

An assistant that knows everything about you and supports you around the clock
is not sold anywhere. You cannot buy one---and, we have argued, you should not
entrust its upbringing to a vendor-owned environment, because the same
accumulation that makes it valuable is what makes
leaving expensive, and defaults are deciding today where that accumulation lives.
This paper offered a working alternative default: an architecture in which the
model authors small declarative applications, a host the user runs executes them,
and everything that accumulates---records, schemas, manuals, history---sits as
plain files in a workspace the user owns. In single-trial runs whose transcripts
and infrastructure ship inside the open-source artifact, the evaluation showed
the architecture trustworthy where it must be (22/23 first-attempt valid schemas;
invalid ones never live), its runtime the source of the capabilities that matter
(each ablation deficit mapped to one removed component), and memory-only
portability moving testimony while leaving capability behind (0 of 4 workflows
survived migration).

The limitations are stated where they arise and bear repeating once: the trials
are single runs pending repetition; the hosted alternative is
discussed rather than observed; and the claim that people will in fact nurture
such an assistant over years is not yet evidence but the question this
architecture was built to make answerable. The substrate for that study now
exists. The upbringing is the user's own.

\bibliographystyle{plainnat}
\bibliography{refs}

\end{document}
